\chapter{Нейрон --- не один прибор}
% полная версия: chapters/18_eetypes.tex

В первой главе мы рассортировали нейроны по тому, какое вещество они выбрасывают. Инженер рассортировал бы их иначе: каким прибором каждый из них работает. И оказывается, что под одинаковым видом прячется целый ящик радиодеталей.

\section*{Ящик радиодеталей}

Дырявое ведро из второй главы --- это накопитель с утечкой: он складывает сигналы за короткое окно. С очень маленьким окном он превращается в детектор совпадений. Нейрон, который отвечает только на перемены и тут же замолкает, --- это фильтр, пропускающий изменения. Нейрон, который работает сам, --- метроном, генератор; если его темп подкручивает вход, это генератор, управляемый напряжением. Нейрон, выдающий пачки щелчков, --- генератор пачек, пишущий «тире». Нейрон, щёлкающий в такт звуковой волне, --- детектор фазы.

\pencilfig{18}{\draw[pen, wash=pcAmber!40] (-0.3,-0.6) rectangle (5.6,0.6);
\draw[pen=pcBlue] (0,0) -- (0.2,0) -- (0.3,0.2) -- (0.45,-0.2) -- (0.6,0.2) -- (0.75,-0.2) -- (0.9,0.2) -- (1.0,0) -- (1.2,0);
\draw[pen=pcRed] (1.6,0) -- (1.95,0); \draw[pen=pcRed, line width=1.2pt] (1.95,-0.3) -- (1.95,0.3); \draw[pen=pcRed, line width=1.2pt] (2.1,-0.3) -- (2.1,0.3); \draw[pen=pcRed] (2.1,0) -- (2.45,0);
\foreach \i in {0,1,2} \draw[pen=pcGreen] (2.9+0.3*\i,0) arc (180:0:0.15);
\draw[pen=pcViolet] plot[domain=4.3:5.4, samples=40] (\x, {0.25*sin(deg(\x*8))});}{Под одинаковым видом прячется целый ящик радиодеталей: сопротивление, конденсатор, катушка, генератор.}

Есть и приборы, о которых мы раньше не говорили. \emph{Резонатор} похож на качели: он лучше всего отвечает на толчки определённого ритма, как качели, которые раскачиваются, только если толкать их вовремя. Внутри нейрона роль пружины играет медленный ток, который сопротивляется изменениям. \emph{Накопитель без утечки} держит значение долго, как конденсатор без утечки, --- так работает сеть, удерживающая положение глаз, когда вы смотрите в сторону. Но для этого обратная связь должна быть подогнана почти идеально, и это дорогая настройка. \emph{Защёлка} --- нейрон, который короткий толчок переводит в долгую работу, пока его не выключат, как выключатель с фиксацией; у нейронов, управляющих мышцами, такой режим включает серотонин.

\section*{Прибор-хамелеон}

Главное открытие этой сортировки: это не разные детали, а разные \emph{режимы} одной и той же. Один и тот же нейрон может быть накопителем днём и генератором пачек ночью, в зависимости от того, что ему говорят радиостанции. Инженер назвал бы мозг перепрограммируемой логической матрицей: схема одна, а назначение деталей меняют настройки.

Что нейроны умеют работать во всех этих режимах, измерено. Как мозг пользуется переключением режимов для вычислений --- по большей части гипотеза.

\fullversion{18}
